% !TEX program = pdflatex
\documentclass[11pt]{article}
\usepackage[T1]{fontenc}
\usepackage{lmodern,amsmath,amssymb,amsthm,microtype,booktabs}
\usepackage[margin=1in]{geometry}
\usepackage[hidelinks]{hyperref}
\hypersetup{pdftitle={When a cubic higher product changes a bounded profile}}
\newtheorem{thm}{Theorem}
\newtheorem{lem}[thm]{Lemma}
\newtheorem{cor}[thm]{Corollary}
\newcommand{\Dfd}{D_{\rm fd}}
\newcommand{\per}{\operatorname{per}}
\newcommand{\thick}{\operatorname{thick}}
\newcommand{\Ext}{\operatorname{Ext}}
\newcommand{\rad}{\operatorname{rad}}
\newcommand{\Hom}{\operatorname{Hom}}
\newcommand{\Cone}{\operatorname{Cone}}
\newcommand{\sg}{\mathcal S}
\newcommand{\UU}{\mathcal U}
\newcommand{\kk}{\Bbbk}
\setlength{\emergencystretch}{1em}
\title{When a cubic higher product changes a bounded profile}
\author{Marouane Lamharzi Alaoui\\
\small Research notes on two-term dg kernels and bounded profiles of universal localisations}
\date{18 November 2026}
\begin{document}
\maketitle
\begin{abstract}
For the one-simple algebra of Example 3.11, both choices of the higher
multiplication have exactly two marked bounded defects ordered by inclusion.
Their nonzero defect categories admit no nonzero additive functor commuting
with suspension between them in either direction. For a connected
four-simple algebra we exhibit a change in the inclusion poset itself.
Keeping the graded cohomology algebra fixed and turning on one cubic
multiplication changes a three-element chain to a two-element chain.
A directed triangle supplies the calculation responsible for the change.
All statements hold over any field.
\end{abstract}

\section{What the poset measures}
Use cochain grading, left dg modules, and a field $k$. For a finite
nonpositive dg algebra $A$, write $\Dfd(A)$ for the modules with finite total
cohomology, $\per A=\thick(A)$, and $q\colon\Dfd(A)\to
\sg(A)=\Dfd(A)/\per A$. A detector $E\in\Dfd(A)$ gives a thick subcategory
\[
 \UU_E=\{X\in\Dfd(A)\mid \textstyle\sum_n\dim_k\Ext_A^n(X,E)<\infty\}.
\]
It contains $\per A$. The manuscript's Corollary 3.10 identifies the
realised marked bounded defects with $\UU_E/\per A$. We compare their
inclusion posets. We also ask whether the triangulated categories at equal
positions in two posets are equivalent. These are distinct questions.

\begin{center}
\begin{tabular}{lcc}
\toprule
Fixed graded cohomology & Formal profile count & Nonformal profile count\\
\midrule
One simple, Example 3.11 & 2 & 2\\
Directed triangle & 2 & 1\\
Triangle with an attached loop vertex & 3 & 2\\
\bottomrule
\end{tabular}
\end{center}
In each row the posets are chains. In the last row both maximal defects
are nonzero. All algebras below are explicit strict two-term dg algebras,
so the manuscript realises them by ordinary universal localisations.

\section{A reduction for finite modules}
All the algebras used here have a split augmentation $A\to S=k^m$ with
nilpotent dg kernel $J$, degree-zero primitive idempotents $e_i$, and
$A=A^{-1}\oplus A^0$. Let $S_i$ denote the corresponding simple left
module and put $S=\bigoplus_i S_i$ as a module. Standard truncation
triangles and composition series in the heart give $\thick(S)=\Dfd(A)$.

\begin{lem}\label{cut}
For each $Y\in\Dfd(A)$, $qY$ is a shift of the image of a two-term
dg module $L$ consisting of elements in $JP'$ for a finite approximation $P'$ built
from shifts of the projectives $Ae_i$. In particular any dg ideal
$I$ with $IJ=0$ kills $L$. If $J^2=0$, $qY$ is a finite direct sum of
shifts of the $qS_i$. In all these examples $\UU_S=\per A$.
\end{lem}
\begin{proof}
Take a minimal semi-free resolution $P$ of $Y$ with cells shifts of
$Ae_i$, bounded above in cell degree, with finitely many cells per degree.
Here is a construction. Working down from the highest cohomology of $Y$,
attach cells mapping to cocycles that span the top cohomology of the
cone of the current approximation. Lift each generator at its own
idempotent. The top cone cohomology is killed and each stage is finite.
The union is a resolution. Modulo $J$, each differential is a matrix over
$S$. Split off its nonzero pivots from the top degree down. A pivot at
vertex $i$ lifts to a unit in $e_iA^0e_i$ because $J$ is nilpotent.
Replacing the upper cell by the differential of the lower one gives a
contractible free pair, whose quotient removes it. In a given total degree
the process stabilises because cells have amplitude one and differentials
raise degree. The resulting $P$ satisfies $dP\subset JP$ (also $dA\subset
J$). Thus $\Hom_A(P,S)$ has zero differential and detects every cell.
Finitely many cells imply perfection, and every perfect source has finite
total Ext against $S$. This gives the final assertion. It also proves that
every Ext group between finite-cohomology objects is finite and sufficiently
negative groups vanish, by building the target from the simples.

For $Y\ne0$ take $a<\inf H(Y)$, and let $P'$ be generated by cells of
degree at least $a$. The differential of a cell uses only cells of higher
degree, so $P'$ is a finite cell submodule. The quotient $Q=P/P'$ is
concentrated in degrees $\le a-1$. The connecting map $Q\to P'[1]$ factors
in the derived category through $(\tau^{\le a}P')[1]$ by the standard
t-structure. The long exact sequence and $H^{\le a}(P)=0$ make this
factor a quasi-isomorphism. Thus $qY\simeq qQ\simeq qL[1]$ for
$L=\tau^{\le a}P'$. It has terms only in degrees $a-1,a$. The lower
term consists of $A^{-1}$-multiples of cells, and $A^{-1}\subset J$.
The upper term consists of cycles, hence boundaries in $P$ and elements
of $JP\cap P'=JP'$. The equality follows from the graded cell basis.
This proves the assertion about $I$. If $J^2=0$, $L$ is a bounded
complex over the semisimple algebra $S$, where it splits into its cohomology.
The zero object satisfies the assertions directly.
\end{proof}

\section{A connected example in which the poset changes}
Read paths from left to right, so $Ae_i$ consists of paths ending at $i$.
Take the oriented triangle
\[
 0\xrightarrow{a}1\xrightarrow{b}2\xrightarrow{c}0 .
\]
For the four-vertex example adjoin $0\xrightarrow{h}3$ and a loop
$\ell$ at 3. In degree zero take the vertex idempotents, these arrows, and
$s\colon1\to0$. In degree $-1$ take $r\colon1\to0$ and $w\colon0\to0$.
Apart from idempotent actions the only nonzero products and differential are
\[
                   bc=s,\qquad ar=w,\qquad dr=s .
\]
These rules give an associative dg algebra $A$: every product of three
radical factors vanishes, and $s$ annihilates the radical on both sides.
The formal comparison $F=H^*(A)$ has the original arrows, the class of
$w$, and radical square zero. Both have the same connected cohomology
quiver. Splitting off the contractible pair $r\mapsto s$, the only
tree term in transfer on arrow triples is $a\,r=w$, from contracting
$bc$. Thus $m_3(a,b,c)=\pm w$, and all other reduced triples vanish.
Idempotents are in the base $S$. Operations of arity at least four vanish
by degree. The value on $(a,b,c)$ of any relative normalized Hochschild
coboundary is zero, since the products of adjacent arrows vanish and
the radical annihilates $kw$. This verifies the nonzero class and also
explains the choice of relations.

\begin{thm}\label{triangle}
For the triangle, $\sg(A)=0$ and $F$ has exactly two profiles.
For the connected four-vertex example, $A$ has exactly two profiles
and $F$ has exactly three. For $F$ in the latter case the intermediate
defect is generated, as a thick subcategory in the quotient, by the
three old vertex simples.
\end{thm}
\begin{proof}
Write $P_i=Ae_i$. In the triangle, right multiplication by $c$ embeds
$P_2$ in $\rad P_0$ with basis images $c,s$. Its quotient has basis
$r,w$ with $ar=w$ and zero differential, hence is $P_1[1]$.
Also $\rad P_1=S_0$ and $\rad P_2=S_1$. The short exact sequences for
these projectives successively make $S_0,S_1,S_2$ perfect. They generate
the triangle's finite-cohomology category.

Adding the bridge and loop leaves $P_0,P_1,P_2$ intact. Set
$e=e_0+e_1+e_2$ and $T=k[\ell]/(\ell^2)$. For any left dg module $X$,
$eX$ is a submodule since $e_3Ae=0$. For $X\in\Dfd(A)$, $eX$ is perfect
by the triangle calculation applied to $eAe$ and by the identification
of its three projectives with $Ae_0,Ae_1,Ae_2$. The quotient $X/eX$
is inflated from $T$. Inflation takes $T$ to a perfect object by the exact
sequence $0\to S_0\xrightarrow{1\mapsto h}P_3\to T\to0$. Over $T$,
Lemma~\ref{cut} reduces finite objects modulo perfects to sums of shifted
copies of its simple. It follows that every $qX$ is a sum of shifts of
$qS_3$. This simple survives: $X\mapsto e_3X$ takes the four projectives
to $0,0,0,T$ and $S_3$ to $k$, whose $T$-resolution repeats multiplication
by $\ell$ indefinitely. Every thick subcategory of the quotient containing
a nonzero object therefore contains $qS_3$ and is the entire quotient.
The detectors $S$ and $0$ give the two endpoints.

For $F$ Lemma~\ref{cut} gives a sum of shifted simples for every object
of the quotient. The radical of each projective is the direct sum of
the incoming arrow simples, with one additional shift for an arrow of
degree $-1$. Thus $qS_i\simeq q(\rad Fe_i)[1]$. Each old simple has its
old predecessor as a shifted direct summand and generates all old simples
in the quotient. They survive: iterating the radical resolution follows
the triangle arbitrarily many times; its differentials vanish against
$S$. For the triangle alone this proves the two endpoints exhaust the
poset. With vertex 3, $qS_3\simeq qS_3[1]\oplus qS_0[1]$ generates all
four. A thick subcategory can consequently only be zero, the one generated
by the old simples, or the whole quotient. The detector $E=S_3$ realises
the middle option. Resolutions of old simples use only old projectives,
so their Ext to $S_3$ vanishes in every degree. Resolving $S_3$ follows
the loop indefinitely and gives nonzero self-Ext in every nonnegative
degree. Again $S$ and $0$ realise the endpoints.
\end{proof}

\paragraph{Independent calculation on the triangle.}
The reduced bar dual over $S$ has the reversed triangle of degree-one
arrows and one degree-two loop $y$ with $dy=p$, where $p$ is the
reversed, dual three-arrow path. Rescale $y$ to fix the sign. The projection
to the path algebra on the three arrows modulo $p$ is a quasi-isomorphism.
Indeed expand each $y$ into $p$ and group basis words by the resulting
path. The occurrences of $p$ in it are disjoint, since the three arrows
are distinct and $p$ has no proper prefix equal to a suffix. In such a
group each occurrence independently contributes the complex
$ky\xrightarrow{d} kp$ with the cochain tensor signs. If there is an
occurrence, the group is contractible. The surviving paths avoiding $p$
have counts $3,3,3,2,1$ in degrees zero through four, respectively;
every path of length five meets $p$. In particular $\dim\Ext_A^*(S,S)=12$.
The projective argument has the same count: $S_0$ has one minimal cell
in each degree $0,-1,-2$ (differentials use $c,b,r$); adjoining $P_1$
and then $P_2$ at the top of the shifted resolutions gives the counts for
$S_1,S_2$. In the formal algebra the entire reduced bar differential
vanishes, and the triangle supplies paths of arbitrary length.

\section{The one-simple comparison in Question 3.21}
For the manuscript's example use its quasi-isomorphic five-dimensional
model $B$. Its radical basis is $t,c$ in degree zero and $r,w$ in degree
$-1$, with $t^2=c$, $tr=w$, $dr=c$ and all remaining radical products
zero. Put $B_f=H^*(B)$ with zero differential. It has square-zero
radical spanned by $t,w$. The earlier continuation note verifies the
quasi-isomorphism with Example 3.11 and the nonzero cubic multiplication.

\begin{thm}\label{one}
For $D=B,B_f$, every object of $\sg(D)$ is a finite direct sum of shifts
of $qk$. Each inclusion poset in Question 3.21 consists of exactly
$0\subsetneq\sg(D)$. In $\sg(B)$ and $\sg(B_f)$ respectively,
\[
            qk\simeq qk[3]\oplus qk[4],
  \qquad    qk\simeq qk[1]\oplus qk[2].
\]
\end{thm}
\begin{proof}
The square-zero case follows from Lemma~\ref{cut}. Its augmentation-ideal
triangle gives the displayed decomposition, since the radical is
$k\oplus k[1]$.

For $B$ put $I=\langle c,r,w\rangle$ and $J=\langle t,c,r,w\rangle$.
Here $IJ=0$, while $JI=kw$. Lemma~\ref{cut} makes $L$ a two-term
complex of finite modules over $G=B/I=k[t]/(t^2)$. We describe every
such complex in $D^b(G)$. Its two cohomology modules are direct sums
of $G$ and $k$, by the square-zero Jordan decomposition. The connecting
class of the truncation triangle belongs to $\Ext_G^2(H^aL,H^{a-1}L)$.
The free summands drop out in either variable since $G$ is Frobenius.
The resolution with differential $t$ gives $\Ext_G^2(k,k)=k$.
The remaining matrix of scalars reduces by independent changes of basis
to an identity block and a zero block. Cones of the resulting blocks
give shifts of $G$, $k$, and the two-term complex $[G\xrightarrow{t}G]$.
For the last block the two-extension
$0\to k\to G\xrightarrow{t}G\to k\to0$ represents a generator, as
seen by the same minimal resolution. This classification includes
arbitrary finite modules as terms; termwise projectivity is unnecessary.

Now restrict to $B$. Inclusion $kw\hookrightarrow I$ is a left dg
module quasi-isomorphism, as $r\mapsto c$ is acyclic in the vector
space quotient. The triangle $I\to B\to G\to I[1]$ gives $qG\simeq qk[2]$.
Right multiplication by $t$ gives an endomorphism of this triangle.
It is zero on $I$ since $IJ=0$. Hence it is zero on $qG$ by naturality.
The cone of $t$ therefore splits in the quotient. This proves the
decomposition for all objects. In fact $J$ itself is the complex
$[G\xrightarrow{t}G]$ in degrees $-1,0$, with generators $r,t$ and
$t r=w,\ t t=c$. Since $qk\simeq qJ[1]$, we get the asserted shifts.

In either case $k$ is nonperfect, by the infinite self-Ext computation
below. Any nonzero sum of shifted copies of $qk$ generates the quotient
as a thick subcategory. More explicitly, for a profile containing a
nonperfect $Y$, passage to $\Dfd(D)/\UU_E$ kills the sum representing
$qY$, hence kills $k$ as a summand. The kernel of the quotient by a
thick subcategory is that subcategory, so $k\in\UU_E$ and $\thick(k)=
\Dfd(D)$. The endpoints are $\UU_S$ and $\UU_0$.
\end{proof}

\section{Growth distinguishes the two nonzero defects}
For the nonformal algebra the reduced bar dual over $k$ is
$K=k\langle x,y\rangle$, $|x|=1$, $|y|=2$, $dx=0$, $dy=x^3$.
This is the dual of $m_3(t,t,t)=w$ on the reduced minimal model;
the sign can be absorbed into $y$. Each degree is finite. In the formal
case the differential vanishes. Write $P_D(s)=\sum_n\dim_k\Ext_D^n(k,k)s^n$.

\begin{lem}\label{series}
Put $z=xy+yx$ and $v=x^2y-yx^2$. A basis for $H^*(K)$ is
$u x^j$, with $u$ any word in $z,v$ and $j\in\{0,1,2\}$. Consequently
\[
 P_B(s)=\frac{1+s+s^2}{1-s^3-s^4},\qquad
 P_{B_f}(s)=\frac{1}{1-s-s^2}.
\]
\end{lem}
\begin{proof}
Both $z,v$ are cycles by the signed Leibniz rule. Classify words by
their initial segment to obtain the graded direct sum
\[
 K=\langle 1,x,x^2\rangle\oplus (yK\oplus x^3K)\oplus zK\oplus vK .
\]
Indeed the original last two summands are $xyK,x^2yK$, and replacing
them by $zK,vK$ adds only terms in $yK$. All displayed blocks are
subcomplexes. The middle block contracts by $h(x^3u)=yu$, $h(yu)=0$;
the differential terms on the suffix cancel in $dh+hd$. Left multiplication
identifies the last blocks with degree shifts of $K$ (with the appropriate
differential signs). Induction on cohomological degree proves the basis.
Counting gives the first series, while counting all words gives the second.
\end{proof}

For $D=B,B_f$ let $\rho_D>1$ be respectively the positive root of
$\rho^4=\rho+1$ and the golden ratio. If $b_n=[s^n]P_D$, then
$b_n\sim c_D\rho_D^n$ for a constant $c_D>0$. Indeed the positive
root of the denominator in $s$ is simple. Equality in the triangle
inequality on its circle forces both powers to have argument zero;
the two exponents are coprime. This gives a unique pole of smallest
modulus and a positive leading coefficient.

\begin{thm}\label{growth}
Every additive functor commuting with suspension from $\sg(B)$ to
$\sg(B_f)$ or in the reverse direction is zero up to isomorphism.
More generally, for either algebra $D$, a finite pair $(X,E)$ has
bounded total Ext precisely when $X\in\per D$ or
$E\in\thick(D_kD)$, where $D_k=\Hom_k(-,k)$. In every other case
$\dim_k\Ext_D^n(X,E)\sim c_{X,E}\rho_D^n$ with $c_{X,E}>0$.
\end{thm}
\begin{proof}
For fixed finite $E$, changing the source by a map with perfect cone
induces isomorphisms on all sufficiently high Ext groups by the long
exact sequence. An isomorphism in the singular quotient is represented
by a roof with both cones perfect. This gives a well-defined sequence
of dimensions up to finitely many terms on each isomorphism class.
Theorem~\ref{one} applied with $E=k$ shows that
\[
 d_D(qX)=\lim_{n\to\infty}\rho_D^{-n}\dim_k\Ext_D^n(X,k)
\]
exists, is finite, and is strictly positive on nonzero objects. It adds
on sums and $d_D(V[j])=\rho_D^{-j}d_D(V)$. If the image of $qk$ under
a suspension-commuting additive functor were nonzero, its source's
splitting identity would force $1=\rho_{\rm target}^{-3}+
\rho_{\rm target}^{-4}$ or $1=\rho_{\rm target}^{-1}+
\rho_{\rm target}^{-2}$ respectively. These identities require different
positive roots greater than one. Hence the image is zero, and every
object is a sum of shifts of the generator.

For the pair statement, the same decomposition holds for the opposite
algebra. For clarity, a strict zigzag for $B^{\rm op}$ and $B$ uses an
algebra $C$ with $C^0=k[t]/(t^4)$ and $C^{-1}=\langle r,tr,rt\rangle$.
The indicated two products with $r$ are independent; higher actions of
$t$ on the bimodule, including $trt$, vanish, and products of negative
elements vanish. Set $dr=t^2$, $d(tr)=d(rt)=t^3$. The Leibniz rule
holds, including on two negative elements since $t^2$ kills the bimodule.
The maps to $B$ and $B^{\rm op}$ sending $t,r$ to the namesake elements
are augmented quasi-isomorphisms: in degree $-1$ the surviving generator
$tr-rt$ maps to $w$ and $-w$. The formal algebra is its own opposite.
Vector space duality on finite-cohomology derived categories sends
$\per(D^{\rm op})$ to $\thick(D_kD)$ and fixes $k$. Hence, modulo this
last subcategory, $E$ is a sum of shifts of $k$, nonempty exactly when
$E$ survives. Since $\operatorname{RHom}_D(X,D_kD)\simeq D_kX$, changing
the target by maps with cones in $\thick(D_kD)$ preserves all high Ext
groups. Combining the two finite decompositions gives a finite sum
of shifted copies of $\Ext_D^*(k,k)$ in high degrees, nonempty exactly
under the stated conditions. Lemma~\ref{series} gives the asymptotic.
\end{proof}

\section{Scope and comparison with prior work}
The cone of the class $z$ in the earlier note cannot give a third
profile. Its nonperfection was established there. Equivalently,
eventual invertibility of $z$ on unbounded $\Ext_B^*(k,E)$ would force
eventual period three of its dimensions, contrary to Theorem~\ref{growth}.
This closes the earlier detector search.

The reduction over the dual numbers is compatible with the classification
of indecomposable complexes by \href{https://arxiv.org/abs/1309.7215}{Amodeo
and Moschetti}. The argument above specifically needs only the two adjacent
cohomology groups. For ordinary radical-square-zero artin algebras,
\href{https://arxiv.org/abs/1104.4006}{Chen} describes the whole singularity
category. Our formal models include a radical generator in degree $-1$.
For commutative dg trivial extensions, \href{https://arxiv.org/abs/1508.00748}
{Avramov, Iyengar, Nasseh and Sather-Wagstaff} prove rigidity of Tor
vanishing. The nonformal one-simple model here is noncommutative and has
$tr=w$ and $rt=0$.

The cubic-relation mechanism has explicit predecessors.
\href{https://webusers.imj-prg.fr/~bernhard.keller/publ/IntroAinfEdinb.pdf}
{Keller} computes a nonzero ternary product for a linearly oriented
four-vertex quiver with a cubic relation.
\href{https://arxiv.org/abs/math/0606144}{Lu, Palmieri, Wu and Zhang}
recover relation coefficients from higher products for connected graded
algebras. The absence of overlaps is classical in
\href{https://doi.org/10.1016/0021-8693(82)90104-1}{Anick's theory of
combinatorially free relations}. The triangle closes the path into a cycle,
which makes forbidding one path bound its length. The attached fourth
vertex preserves a singular component while the other one disappears.
The profile computation also uses the manuscript's realisation theorem.
The notion of far-away orthogonality is studied by
\href{https://arxiv.org/abs/2512.18417}{Kostas, Psaroudakis and
Vit\'oria}. These comparisons identify prior ingredients; priority for
the combined profile results requires a wider review.

\paragraph{Remaining local question.}
The next comparison asks whether two proper two-term dg algebras with
the same graded cohomology and a single simple module can have nonisomorphic
profile posets while both singularity categories survive. Theorem~\ref{one} settles the
one-loop test in the manuscript. The multivertex examples do not decide
this local question. Even the ordinary local algebra
$R=k[u,v]/(u^2,v^2)$ has an intermediate profile: take
$X=R/(u), E=R/(v)$. Multiplication by $u$ gives a periodic minimal
resolution of $X$, and is exact in positive cochain degrees after applying
$\Hom_R(-,E)$. Thus $X$ is nonperfect and lies in $\UU_E$. Tensoring
the two periodic residue-field resolutions over $k$ shows
$\Ext_R^n(k,E)\cong k$ for every $n\ge0$ (the $u$ factor contributes
only its socle in degree zero). A promising local test would have to
alter such a detectable vanishing pair while keeping the entire graded
cohomology fixed. The companion note \emph{A cubic higher multiplication creates infinitely
many bounded profiles} resolves this question using two degree-zero
radical generators and a six-dimensional commutative dg algebra.

\end{document}
